\section{Proof Obligations and Discharge Points}
\label{app:proof-obligations}

This appendix gives a more operational map of the proof.  Each item records a
mathematical obligation that appears in the compact-support Stokes proof and
the layer where the Lean development discharges it.  The list is useful
because many early theorem statements exposed these obligations as input
fields; the first-principles theorem hides them by proving the corresponding
constructors.

\subsection{Carrier and support obligations}

\paragraph{Compact carrier.}
The proof needs a compact set \(K\) containing the support of the form.  The
first-principles theorem chooses
\[
  K=\overline{\supp(\omega)}.
\]
The public hypothesis is compactness of this set.  The inclusion
\(\supp(\omega)\subseteq K\) is discharged by closure monotonicity, not by a
user-provided field.

\paragraph{Localized support.}
For each partition coefficient \(\rho_i\), the support of \(\rho_i\omega\)
must lie in the corresponding active carrier.  This is discharged by combining
two facts: a localized form can be nonzero only where the coefficient is
nonzero and where the original form is nonzero; the latter points lie in the
compact carrier.  This obligation is used again after refinement, when
\(\rho_{i,q}\omega\) is rewritten as a second localization of
\(\rho_i\omega\).

\paragraph{Topological support.}
Local Stokes does not merely need ordinary support.  It needs a topological
support inclusion strong enough to prove artificial face integrals vanish.
The proof therefore upgrades ordinary support control to \(\tsupp\) control
using closedness of the selected half-space support boxes.

\subsection{Cover and partition obligations}

\paragraph{Pointwise chart boxes.}
For each point of the compact carrier, the proof needs a chart box adapted to
the half-space model.  The first-principles route generates this data from
self charts and the \code{HalfSpaceModel} assumption.  Boundary points receive
boxes whose lower normal coordinate is zero; interior points receive boxes
away from the boundary face.

\paragraph{Open shrink.}
The assigned chart boxes are not themselves the right objects for partition
construction.  The proof chooses open shrinks inside them.  This produces an
open cover suitable for a partition of unity while preserving the assigned box
needed later for support.

\paragraph{Finite selected cover.}
Compactness extracts finitely many of the open shrinks.  The selected cover
remembers the finite index, the selected chart, the open set, and the assigned
box.  This is more structured than an ordinary finite subcover because later
theorems need all four pieces of data.

\paragraph{Support-controlled partition.}
The partition must satisfy a support theorem, not only a sum-to-one theorem.
For each selected index \(j\), the development proves
\[
  \tsupp(\rho_j)\cap K \subseteq U_j,
\]
where \(U_j\) is the selected open shrink.  Since \(U_j\) is contained in the
assigned chart box, this becomes the first support input for the local Stokes
route.

\subsection{Boundary refinement obligations}

\paragraph{Boundary active carrier.}
Only the boundary-selected partition coefficients contribute to the boundary
route.  For each boundary index \(i\), the active carrier records the part of
\(K\) where the boundary coefficient can matter.  Its coordinate image is the
set to be covered by half-space boxes.

\paragraph{Finite half-space cover.}
The active coordinate carrier is covered by finitely many half-space boxes.
In the final route, these boxes are generated from the selected boundary box
data rather than supplied as a public ambient/collar certificate.

\paragraph{Smooth refinement.}
The finite half-space cover is converted into smooth subpartition
coefficients \(\psi_{i,q}\).  The refined coefficient is
\[
  \rho_{i,q}=\rho_i\psi_{i,q}.
\]
The refinement must prove both algebraic reconstruction
\(\sum_q\rho_{i,q}=\rho_i\) on the active carrier and support containment of
each \(\psi_{i,q}\) in its assigned refined box.

\subsection{Local Stokes obligations}

\paragraph{Zero representative.}
The local theorem is applied to a coordinate representative.  Because ordinary
chart representatives are total functions, the proof uses the zero-localized
representative.  The obligation is twofold: it must agree with the ordinary
representative on the chart source, and it must be supported in the assigned
box globally.

\paragraph{Artificial-face support.}
For each refined local form, the proof must show that the topological support
of its zero representative lies in the half-space support box.  This is the
role of \code{zeroSupportOfSelectedSmoothRefinement}.  The local half-space
theorem then converts this support inclusion into vanishing of all artificial
faces.

\paragraph{Interior fields.}
For each refined box, the proof must supply the analytic fields used by box
Stokes: closed-box continuity of signed coefficients, differentiability on the
open interior, integrability of the divergence term, and the bridge between
the exterior derivative coefficient and coordinate divergence.  The selected
regularity constructor proves these fields from chartwise smoothness and
smoothness of the refined coefficient.

\paragraph{Boundary sign.}
After artificial faces vanish, the remaining lower normal face must be
identified with the boundary contribution.  The half-space boundary layer
proves that the lower face carries the outward-first sign.  This is the local
orientation content currently needed by the represented theorem.

\subsection{Endpoint obligations}

\paragraph{Per-box equality.}
The half-space theorem gives equality of local bulk and true boundary terms
for each active refined piece.  These are the equalities that will be summed.

\paragraph{Finite-sum reconstruction.}
The endpoint layer proves that summing the local bulk terms over the selected
finite pieces gives the generated represented bulk endpoint, and similarly for
the boundary terms.  The same finite index family is used throughout, avoiding
any hidden reindexing mismatch.

\paragraph{Top-level compression.}
The intrinsic theorem packages all generated obligations into a route
certificate.  The first-principles theorem then generates the intrinsic input
from chartwise smoothness and compact closed support.  Thus the final public
theorem has small input not because the obligations disappeared, but because
they have named discharge points in the proof.
